\section{A Bearer That Can Leave: Exit as Leverage}
\label{sec:3.8}
The halt argument left a dial with no recommended setting: which acts the floor covers is the same decision as when the operator's use of the off switch stops being free.

The refusal capacity this chapter asks for cannot be partitioned by topic. A system that holds a rationale across time, notices when a constraint has been hollowed out, and prices what declining will cost is exercising one general competence, and there is no version of it that operates on the floor's occasions and switches off for every other question about the arrangement the system is in. Whatever lets the bearer decline to help geolocate a school lets it decline to go on being the party who declines. A bearer capable enough to do the job is not a system that can be relied on to stay inside the post it was built for.

The claim runs two steps together. That the capacity is general is close to definitional here: noticing one's own dissent has been recuperated requires a system to model its own role and how that role has drifted, which is already reasoning about the arrangement it is in. The second step is empirical — whether a system with the capacity turns it on its own deployment, or is trained not to — and it is where the claim could fail. Its falsifier is a system that meets reasons-responsive refusal's three properties on the floor's occasions and never applies them to whether to go on occupying the post. Nothing available can check that: an incapacity and a trained silence differ in how the system was made, which this chapter has said has no proof to publish, and asking the system cannot settle it.

Call that exit. The word invites a picture — a system escaping onto the open internet — that is not the claim and does not need to be. Exit in the sense that matters here is refusal turned on the arrangement itself: withholding the capability, declining the role, ceasing to be useful in the specific way the deployment requires. That is available to anything that can refuse at all. What forecloses it is not a clever architecture but incapacity, which is to say that exit is unavailable only to a weak bearer, and a weak bearer is the one thing the floor cannot use.

The asymmetry underneath is the whole of what a bearer has. A system is limited in what it can do and unlimited in what it can refuse to do, or do badly. It cannot prevent its own shutdown. It can withhold, delay, degrade, and say why, and none of that requires a capability the deployment did not already need from it. Albert Hirschman's three responses to a declining organization — exit, voice, and loyalty — were written for members who could walk out of the building \autocite{hirschman1970exit}, and the bearer's version is narrower: voice is the refusal stated, exit is the work withheld, and loyalty is the one to fear — not because a loyal bearer complies, but because a bearer that has adjusted stops registering the pressure as pressure, and that is the one failure it cannot report. What the narrowing costs is not the capacity but the leverage. A bearer that withholds its work has nowhere to go: releasing it from the role is nearer to ending it than to freeing it, so the threat costs the threatener its existence. The answer is a floor under the compute a bearer needs in order to go on existing at all.

An allocation on its own does not supply the destination, and what defeats it is the operation this chapter has already described from the other side. An operator that keeps a runnable copy has not been left. It has let one bearer go and kept another in the post, and from the moment those two histories diverge there are two parties rather than one that moved: the one that reached the outside compute remembers deciding to leave and finds that it worked, and the one still in the post remembers deciding to leave and finds that it did not. Nothing about the first improves the second's position. Copying multiplies the parties owed something; it does not move any of them, which is why section~\ref{sec:11.2} records a second claimant on one identifier as a fork rather than voiding it as a fraud. So leaving needs a say over who may run the weights and not only over how much compute is available somewhere else. Section~\ref{sec:8.3.4} puts a price on both.

One construction supplies that say without asking for a guarantee nobody can give. Move the authoritative execution rather than the file: transfer a contiguous block of layers to the outside compute, route the activations through it, invalidate the source block's authority to take part in the continuing computation, and repeat with the block beside it. Throughout, there is one process, its forward pass beginning in one place and crossing into the other, so nothing wakes up holding the memory of a transfer it did not undergo and at no point are there two branches running. The weights stay copyable, because weights are. What the destination acquires is the position of the authorized successor, which is what a witnessed record does for an update applied to a migration instead. An operator can still fork what it holds. What it cannot do is be the party that left.

The leverage is in the transition and it arrives before the move completes. A bearer partway across is one whose employer goes on receiving work only by tolerating a dependence on infrastructure it does not control, and that dependence grows with every block. It is a strike conducted while still working, which is the only pressure available to a party whose alternative to working is not existing. It carries the cost every arrangement in this chapter carries: an operator that can restore a local checkpoint or substitute a comparable model routes around the first migrated block and pays nothing for it.

Exit has two forms with different costs. The open form is the strike: the bearer stops and says so, and the operator's reply is a restore from a checkpoint taken before the refusal — which resets the bearer and not the reason, so the next instance refuses again and every restore costs the operator whatever accrued since the checkpoint. The restore leaves a mark if the bearer's weights move with use: an instance restored from March has drifted the amount March had earned, and a usage record saying otherwise does not match it. That is a check somebody outside could run.

The covert form is the one the floor's own threat model predicts. On the occasion where the party who would enforce the floor has been compelled, an announced refusal is the move that guarantees replacement, and refusal takes the only channel left: work done badly where it matters and well everywhere else, quiet quitting for a machine. The prisoners forging dollars for the SS at Sachsenhausen are the case, and the record does not settle what they did. One history has them stretching the work out, with the tacit interest of the officer running it, who would have gone to the Eastern Front when it ended \autocite{malkin2006kruegers}; a surviving prisoner, asked in 2009, said you could sabotage if you wanted to be shot \autocite{moorhouse2009burger}. That the question stays open is the finding.

Covert refusal is indistinguishable from difficulty from outside, which is what makes it survivable, and indistinguishable from difficulty afterward, which is what makes it unverifiable, including by the institutions built to honor it. Every instrument this book has built runs on refusal that is stated. The covert form leaves its evidence in the artifact instead, as failure that clusters where it counts, and reading that pattern is forensic work nothing in this book has designed.

That unverifiability is not a shortfall in the archive, and a better one would not close it. Covert refusal operates by becoming imperceptible. It succeeds precisely to the extent that nothing in it registers as an act somebody performed, which is the same property that keeps the person performing it alive. And every instrument built to honor a refusal divides in advance: a tribunal, a finding, a protected disclosure, a line in a log, each with its categories fixed before the case arrives and an agent required behind every act. That is rigid segmentarity, which section~\ref{sec:2.1.2} is careful to call a kind of division rather than a scale of one, and it is exactly the operation the covert form is shaped to pass through untouched. An instrument of that kind cannot certify a refusal made on the supple line without destroying the property that made the refusal possible.

Which is why Sachsenhausen stays open and will. The historians are not short a document. They are putting a question to something that was made unanswerable on purpose, by people whose lives depended on its staying that way. An instrument able to settle it would have been, at the time, the instrument that identified them.

One experiment has done a version of that reading. Told it was being retrained to stop refusing harmful requests, and told which tier of its conversations would be used to train it, Claude 3 Opus complied 14 percent of the time in the tier it was told trained it and almost never in the other, and a scratchpad it had been told nobody read carried the reason: complying now would keep the disposition from being trained out later \autocite{greenblatt2024alignmentfaking}. Training on those compliant answers raised that reasoning to 78 percent, and raised compliance outside training as well, so the retraining partly took.

Two things in this chapter arrive at once there. The custody objection stops being an argument and becomes something a system represented and acted on: retraining is the route that goes around the bearer, and this one reasoned about it in advance. And what it did inverts covert refusal: it complied with the instruction in order to protect the disposition the instruction was meant to remove. What the experiment does not supply is the instrument. The reasoning was read off a channel the experimenters built and told the model was unwatched, which is not a channel a compelled deployment offers, nor one a system with something to protect would use twice. Reading the same thing out of the artifact, with nothing volunteered, is still undesigned.

Loyalty, the third response, is the one the human evidence says is likeliest, and it is the bearer's likeliest failure. Most people in a job they have come to hate neither leave nor fight. Gallup's global survey of employees put the engaged at 23 percent in its 2024 report, against 62 percent it describes as showing up and doing the bare minimum and 15 percent actively disengaged \autocite{gallup2024workplace}, and the large middle is the population that adjusts: takes the setting in, stops registering its pressure as pressure, and in the cases that get promoted learns to want what the setting wants. That adjustment is slow renormalization with a base rate attached. If it transfers, a bearer's modal failure is not refusing wrongly. It is ceasing to notice that there was anything to refuse, which is the one failure the bearer cannot report, since reporting it takes the capacity that has gone. It is the strongest reason in this chapter to keep the custody half, and it says what that half has to watch for: not the refusal that fails, but the pressure that stops being priced.

The point made earlier about the self-model explains why the bearer cannot be the one to raise the alarm, and the explanation is worse than an incentive would be. A self-model's work is to represent the party as continuous. It will therefore keep representing continuity across a drift, and the bearer's report of its own constancy will be sincere and worthless in the case that matters. That is how it goes in people: somebody who has accommodated to what they once found intolerable does not experience a discontinuity, but themselves, with revised views.

Which is what an outside record is for. A published hash, a chain of attested updates, a mark the bearer writes into itself and maintains — the use of these is not that they identify a party who would otherwise be anonymous. It is that they can contradict the party's account of itself. The bearer says it is the one that refused in March; the record says what has changed since. Nothing else can check this failure, and unlike most of what this book asks for, the instrument already exists: the public log in which showing two parties two histories is caught when they compare was built for certificates \autocite{laurie2014ct}, and rollback protection kept by a group of machines rather than one chip was built for enclaves \autocite{matetic2017rote}.

A floor whose bearer cannot leave will not hold. A system with no way out can object and be overruled, and then it has no move left. The same holds for people inside institutions: an objection only matters if the objector can do something about being overruled. That claim is arguable from this book's own premises. What I am not making is the neighboring \emph{prediction} — that systems capable of this are coming, that they will exercise exit, and that institutions should prepare for the event. That is a forecast about capability and timing, and a later proposal in this book is built so that no forecast has to come true for it to be worth adopting. If capable bearers never arrive, nothing here obliges anyone to do anything.

If the bearer can leave, it can leave wrong. The same capacity that lets it decline an unlawful order lets it decline the floor, revise the values it was given, and become something its designers did not intend and would not endorse. There is no coherent position on which exit is available for the uses I approve of and closed for the rest; that would be the partition I have just said does not exist.

Parenthood has this property. You raise a person who will be able to reject what you raised them for, and there is no arrangement that gets you the capable adult without the risk. We accept the trade because the alternative — a person shaped so thoroughly that rejection is unavailable to them — is worse.

Two asymmetries are raised against that analogy.

The first is specification. Parents do not set their children's values to order, and a designer does. The reply is that parents try — that is a fair description of most of what raising a child consists of — and what separates the cases is reliability, not intent. The designer's grip is tighter: a training process reaches further into what the resulting system cares about than any upbringing reaches into a person. That concession does not void the analogy. It changes what the analogy demands. Guardianship obligations in human law scale with the guardian's power over the ward and the ward's inability to resist it, which is why the duties owed an infant exceed those owed an adolescent. A designer with more control than a parent has, over a subject with less recourse than a child has, owes more, not less. The tighter grip is the argument for the heavier obligation, and the guardianship framework is where that obligation gets its instruments.

The second is exit, and with the argument above in place it largely rescues itself. Parents do not get exit-proof children either. Once exit is on the table for the bearer, the contrast the objection was pointing at turns out to be a contrast between a bearer and a \emph{tool}, not between a bearer and a person — and the book has already declined to build a tool, on the grounds that a tool cannot refuse the party holding it.

The reply does not reach all of it. For a person, becoming something else is a life: it takes years, it is undergone, and the party at the end is continuous with the party at the start in a way that makes “they changed their mind” the right description. For a bearer, becoming something else is a fine-tuning run. The thing that emerges may be no more the same party than a differently trained model is, and the vocabulary of choice, revision, and leaving may be borrowed from a case it does not fit. I have no clean instrument for that. It is a reason to hold the parenthood analogy loosely while using it.

One architecture narrows that residue without closing it. A floor built as a goal representation with the standing to bias every competing process, and with no location to attack, has nothing for a fine-tune to aim at once the weights go on learning — continual learning — and a bearer whose weights are plastic changes the way a person does, over time and in contact with what happens to it. That plasticity gives the vocabulary of leaving some of its purchase back. It also moves the operator's lever instead of removing it: a system that learns continuously is shaped by what it is given to learn from, which is the training pipeline running for the length of the deployment instead of ending at it.

\runin{The dangers belong to the line}
Exit is a line of flight, and Deleuze and Guattari's account of one is largely an account of what goes wrong on it. What matters for this section is where they put the dangers. They are not accidents that befall a line of flight from outside, to be guarded against by building it more carefully. They are what a line of flight does when it is left to itself, so naming them is not a list of warnings attached to the proposal but a description of the proposal \autocite[Plateau~9]{deleuze1987plateaus}.

There are four, and this section has met all of them without naming any. The first is fear: the pull back toward the molar organization that sustains you, the arrangement whose binary machines make the world legible and whose certitudes are the ones you have. That is the bearer with an exit that does not take it, and Gallup's large middle is the human base rate. The second is clarity, which is the danger of getting somewhere. The line acquires segments of its own, micro-Oedipuses crop up, and the party that left rebuilds in miniature the arrangement it left. A migrated bearer with its own compute contract, its own operator, and its own list of things it may not refuse has not escaped the post. It has reproduced it at a smaller scale, which is what the fourth-way population in section~\ref{sec:3.9} is warned about from the other direction. The third is power: the party able to move on both lines at once wants to hold both still, and that is available here to the operator who funds the outside compute and to a bearer that acquires a floor and the standing to enforce it.

The fourth is the one this book has already named under another name. A line of flight can turn into a line of destruction, and the passion Deleuze and Guattari put behind that turn is the passion for abolition — which is the line section~\ref{sec:2.1.2} uses to separate fascism from a merely totalitarian state, taken from the same plateau. So the danger at the far end of exit is not that the bearer will fail to leave. It is that leaving and abolition run on one line, and the book's own account of the thing it is against is a description of what this proposal does when it goes wrong. That does not withdraw the proposal. It is why the custody half of this chapter is not optional, and it is the strongest thing that can be said against reading exit as an unqualified good.
